\clearpage
\VerdierPhysicalPageAnchor{0043}%
\VerdierSemanticAnchor{ch1-1-3-8-composition-cont}
与 $G^\circ$-范畴 $(\mathcal C,G^\circ,T,c)$ 相伴的分次范畴
$\mathcal C[G]$ 中的复合定义 (1.3.2.5) 相容。

\VerdierSemanticAnchor{ch1-1-3-9}
\noindent\textbf{1.3.9.}\quad
设 $(\mathcal C,G,T,c)$ 为一个 $G$-范畴。$(\mathcal C,G,T,c)$ 的反范畴是
下述 $G^\circ$-范畴（$G^\circ$ 是 $G$ 的反群）
$(\mathcal C^\circ,G^\circ,T^\circ,c^\circ)$。其中 $\mathcal C^\circ$ 是
$\mathcal C$ 的反范畴。对 $g\in G^\circ$，函子 $T^\circ(g)$ 定义为

\VerdierSemanticAnchor{ch1-1-3-9-eq1}
\begin{equation}
  \begin{aligned}
    T^\circ(g)X &= T(g^{-1})X,
      & X&\in\operatorname{Ob}(\mathcal C^\circ)
             =\operatorname{Ob}(\mathcal C),\\
    T^\circ(g)u &= T(g^{-1})u,
      & u&\in\operatorname{Fl}(\mathcal C^\circ)
             =\operatorname{Fl}(\mathcal C),
  \end{aligned}
  \tag{1.3.9.1}
\end{equation}
而 $\mathcal C^\circ$ 中的同构 $c^\circ(g,h)$ 定义为

\VerdierSemanticAnchor{ch1-1-3-9-eq2}
\begin{equation}
  \begin{aligned}
    c^\circ(g,h)&\colon
      T^\circ(h\perp g)\longrightarrow T^\circ(h)T^\circ(g),\\
    c^\circ(g,h)&=c^{-1}(g^{-1},h^{-1}).
  \end{aligned}
  \tag{1.3.9.2}
\end{equation}
（我们以 $\perp$ 表示 $G^\circ$ 中的复合。）

\VerdierSemanticAnchor{ch1-1-3-10}
\noindent\textbf{注 1.3.10.}\quad
设 $\mathcal C[G^\circ]\longrightarrow G^\circ$ 为与范畴
$(\mathcal C,G,T,c)$ 相伴的 $G^\circ$ 型分次范畴。该范畴带有一个典范劈分。
再令 $\mathcal C[G^\circ]^\circ\longrightarrow G$ 为取反范畴所得的余纤维化范畴。
它带有一个典范余劈分；又因为 $G$ 是群，它也纤维化于 $G$ 之上。因此，
$\mathcal C[G^\circ]^\circ\longrightarrow G$ 带有一个典范劈分，从而给出一个
典范同构于 $(\mathcal C,G,T,c)$ 的反范畴的 $G^\circ$-范畴；这一点可立即验证。

\VerdierSemanticAnchor{ch1-1-3-11}
\noindent\textbf{1.3.11.}\quad
设 $(\mathcal C',G',T',c')$ 为一个 $G'$-范畴，
$(\mathcal C'',G'',T'',c'')$ 为一个 $G''$-范畴。以显然的方式，范畴
$\mathcal C'\times\mathcal C''$ 典范地带有一个 $G'\times G''$-范畴结构。

\VerdierSemanticAnchor{ch1-1-3-11-product}
所得 $G'\times G''$-范畴称为 $G'$-范畴 $(\mathcal C',G',T',c')$ 与
$G''$-范畴 $(\mathcal C'',G'',T'',c'')$ 的\emph{积}。

\VerdierSemanticAnchor{ch1-1-3-12}
\noindent\textbf{1.3.12.}\quad
设 $(\mathcal C,G^\circ,T,c)$ 为一个 $G^\circ$-范畴。在实际中，我们通常对
$G^\circ$ 在 $\mathcal C$ 上的作用采用右侧记法。对 $\mathcal C$ 的每个对象
$X$，置

\VerdierSemanticAnchor{ch1-1-3-12-eq1}
\begin{equation}
  T(g)X=X[g]
  \tag{1.3.12.1}
\end{equation}
同样，对 $\mathcal C$ 中每个态射 $u$，置

\VerdierSemanticAnchor{ch1-1-3-12-eq2}
\begin{equation}
  T(g)u=u[g].
  \tag{1.3.12.2}
\end{equation}
于是，对每个对象 $X$，同构 $c(g,h)$ 定义同构

\VerdierSemanticAnchor{ch1-1-3-12-eq3}
\begin{equation}
  c(g,h)(X)\colon X[gh]\longrightarrow X[g][h].
  \tag{1.3.12.3}
\end{equation}

\VerdierSemanticAnchor{ch1-1-3-12-notation}
这一记法尤其有一个优点：它把 $G^\circ$ 在 $\mathcal C$ 上的作用变成 $G$ 在
$\mathcal C$ 上的右作用。
